\chapter{Apr.~30 --- Hodge Theory}

\section{Hodge Decomposition Theorem}

\begin{remark}
  For this lecture, we will always denote
  Cech cohomology by $\widecheck{H}^i$.
  We want to see how this is related
  to the singular cohomology
  $H^i_{\mathrm{sing}}(X^{\mathrm{an}}, \C)$.
\end{remark}

\begin{theorem}[Hodge decomposition theorem]
  For $X \subseteq \PP^n_\C$ a smooth
  projective variety of dimension $n$, there
  exists a decomposition
  \[
    H^k_{\mathrm{sing}}(X^{\mathrm{an}}, \C)
    = \bigoplus_{p + q = k} H^{p, q}
  \]
  for $0 \le k \le 2n$. Furthermore,
  \begin{enumerate}
    \item $\overline{H^{p, q}} = H^{q, p}$ (where we view $H^k_{\mathrm{sing}}(X^{\mathrm{an}}, \C) = H^k_{\mathrm{sing}}(X^{\mathrm{an}}, \R)$ by the universal coefficients theorem and $\overline{H^{p, q}}$ denotes complex conjugation);
    \item $H^{p, q} = H^q(X, \Omega_X^p)$
      ($= H^q(X^{\mathrm{an}}, \Omega_{X}^{p, \mathrm{an}})$ by GAGA);
    \item writing $h^{p, q} = \dim H^{p, q}$,
      we have $h^{p, q} = h^{q, p}$
      (by (1), also called \emph{Hodge duality})
      and $h^{p, q} = h^{n - p, n - q}$
      (by Serre duality).
  \end{enumerate}
\end{theorem}

\begin{remark}
  Note that $b_i = \dim H^i_{\mathrm{sing}}(X^{\mathrm{an}}, \C) = \sum_{p + q = i} h^{p, q}$,
  so the Betti numbers
  $b_i$ are even for $i$ odd.
\end{remark}

\begin{example}
  Let $C$ be a genus $g$ curve. Then
  for $k = 1$, we have
  \[
    h^{1, 1} = h^{0, 0}
    = h^0(C, \OO_C) = 1.
  \]
  Note that we have the diamond
  \[
    \begin{tikzcd}[row sep=small, column sep=tiny]
      & h^{1, 1} = 1 \\
      h^{0, 1} = g & & h^{1, 0} = g \\
      & h^{0, 0} = 1
    \end{tikzcd}
  \]
\end{example}

\begin{remark}
  We want to explain where the $H^{p, q}$
  come from and why it is related to
  $H^q(X, \Omega_X^p)$.
\end{remark}

\section{Cohomology of Sheaves}

\begin{remark}
  Let $X$ be a topological space and
  $\mathcal{F}$ a sheaf of abelian groups
  (e.g. $X$ a complex manifold with
  $\mathcal{F} = \underline{\C}$ or $\OO_X$).
  Recall that we have defined Cech
  cohomology, which we now denote by
  $\widecheck{H}^i(X, \mathcal{F})$.
\end{remark}

\begin{definition}
  A sheaf $\mathcal{F}$ is \emph{flasque}
  if $\mathcal{F}(U) \to \mathcal{F}(V)$
  is surjective for all $V \subseteq U \subseteq X$
  open.
\end{definition}


\begin{example}
  Let $\mathcal{G}$ be a sheaf. Define
  \emph{Godement sheaf} $\mathrm{God}(\mathcal{G})$
  by
  \[
    U \longmapsto
    \mathrm{God}(\mathcal{G})(U)
    := \prod_{x \in U} \mathcal{G}_x.
  \]
  Then we can observe that
  $\mathrm{God}(\mathcal{G})$
  is flasque.
\end{example}

\begin{definition}
  A \emph{flasque resolution} of a sheaf
  $\mathcal{F}$ is a complex of sheaves
  \[
    \begin{tikzcd}
      0 \ar[r] & \mathcal{F}
      \ar[r] & \mathcal{I}^0
      \ar[r, "d_0"] & \mathcal{I}^1
      \ar[r, "d_1"] & \cdots
    \end{tikzcd}
  \]
  such that
  \begin{enumerate}
    \item $d^2 = 0$;
    \item $\im d_{i - 1} = \ker d_i$ for
      all $i$;
    \item each $\mathcal{I}^i$ is flasque.
  \end{enumerate}
\end{definition}

\begin{example}[Godement resolution]
  Let $\mathcal{F}$ be a sheaf. Then
  \[
    \begin{tikzcd}
      0 \ar[r] & \mathcal{F}
      \ar[r, "\alpha"] & \mathrm{God}(\mathcal{F}) \ar[d, dashed] \ar[r, "d_0"] & \mathrm{God}(\coker \alpha) \ar[d, dashed] \ar[r] & \cdots \\
                       & & \coker \alpha \ar[ur, dashed] & \coker d_0 \ar[ur, dashed]
    \end{tikzcd}
  \]
  is a (canonical) flasque resolution of
  $\mathcal{F}$.
\end{example}

\begin{definition}[Sheaf cohomology]
  Let $X$ be a topological space and
  $\mathcal{F}$ a sheaf of abelian groups.
  Then we can define
  \[
    H^i(X, \mathcal{F})
    := H^i(\Gamma(X, \mathcal{I}^\bullet)),
  \]
  where $\mathcal{F} \to \mathcal{I}^\bullet$ is a flasque resolution of $\mathcal{F}$ (e.g. the Godement resolution).
  One can check that this definition is
  independent of the choice of flasque resolution.
\end{definition}

\begin{example}
  If $X$ is locally contractible
  (e.g. a manifold), then
  \[
    H^k(X, \underline{\R})
    \cong H^k_{\mathrm{sing}}(X, \R).
  \]
  To see this, consider the complex of sheaves
  \[
    \begin{tikzcd}
      0 \ar[r] & \underline{\R}
      \ar[r] & \mathcal{C}^0_{\mathrm{sing}}
      \ar[r, "\partial"] & \mathcal{C}^1_{\mathrm{sing}}
      \ar[r, "\partial"] & \cdots
    \end{tikzcd}
  \]
  where $\mathcal{C}^p_{\mathrm{sing}}$ is
  the sheaf defined by
  $U \mapsto C^p_{\mathrm{sing}}(U, \R)$.
  We can note that:
  \begin{enumerate}
    \item $\partial^2 = 0$ (from algebraic topology);
    \item the sequence is exact
      (since $X$ is locally contractible);
    \item each $\mathcal{C}^p_{\mathrm{sing}}$ is flasque (essentially by the definition of a cochain).
  \end{enumerate}
  Thus this complex is a flasque resolution,
  so we
  see that
  \[
    H^k(X, \underline{\R}) = \frac{\ker(\partial : C^k_{\mathrm{sing}}(X, \R) \to C^{k + 1}_{\mathrm{sing}}(X, \R))}{{\im(\partial : C^{k - 1}_{\mathrm{sing}}(X, \R) \to C^k_{\mathrm{sing}}(X, \R))}}
    = H^k_{\mathrm{sing}}(X, \R).
  \]
  In fact, this works more generally for
  any ring $R$, not just $\R$.
\end{example}

\pagebreak

\begin{theorem}
  If $X$ is Hausdorff and paracompact,
  then
  \[
    \widecheck{H}^i(X, \mathcal{F})
    \longrightarrow H^i(X, \mathcal{F})
  \]
  is an isomorphism.
\end{theorem}

\begin{example}
  Let $X$ be a smooth manifold. Consider
  the de Rham complex
  \[
    \begin{tikzcd}
      0 \ar[r] & \underline{\R}
      \ar[r] & \mathcal{A}^0
      \ar[r] & \mathcal{A}^1
      \ar[r] & \cdots
    \end{tikzcd}
  \]
  where $\mathcal{A}^p$ is the sheaf
  defined by
  $U \mapsto \mathcal{A}^p(U)$ (the space of
  smooth $p$-forms on $U$). We have:
  \begin{enumerate}
    \item exactness: this is by the Poincar\'e lemma;
    \item $d^2 = 0$: this follows from differentiation rules;
    \item flasqueness: this fails, even
      for $\mathcal{A}^0$.
  \end{enumerate}
  However, $\mathcal{A}^p$ is still
  a \emph{fine} sheaf, so by Theorem
  \ref{thm:fine} below, we still have
  \[
    H^p(X, \underline{\R})
    = \frac{\ker(d : \mathcal{A}^p(X) \to \mathcal{A}^{p + 1}(X))}{\im(d : \mathcal{A}^{p - 1}(X) \to \mathcal{A}^p(X))}
    = H^p_{\mathrm{dR}}(X, \R).
  \]
\end{example}

\begin{definition}
  A sheaf $\mathcal{F}$ is \emph{fine} if
  it can be viewed as a sheaf of
  $\mathcal{A}$-modules for some sheaf of
  rings $\mathcal{A}$ on $X$ such that
  for every open cover $\{U_i\}_{i \in I}$
  of $X$, there exist $f_i \in \mathcal{A}(U_i)$
  such that
  $\sum_{i \in I} (f_i)_x = 1$.
\end{definition}

\begin{example}
  Each $\mathcal{A}^p$ is fine, as we can
  view $\mathcal{A}^p$
  as a sheaf of $\mathcal{A}$-modules,
  where $\mathcal{A}$ is the sheaf of
  smooth functions. Note that
  $\mathcal{A}$ admits partitions of unity.
\end{example}

\begin{theorem}\label{thm:fine}
  If $\mathcal{F}$ is a sheaf and admits a
  resolution
  $0 \to \mathcal{F} \to \mathcal{I}^\bullet$
  by fine sheaves, then
  \[
    H^i(X, \mathcal{F})
    = H^i(\Gamma(X, \mathcal{I}^\bullet)).
  \]
\end{theorem}

\section{Complex Geometry}

\begin{remark}
  Let $X$ be a complex manifold
  of complex dimension $n$
  and $\mathcal{A}^k_\C$ the sheaf
  of complex $k$-forms.
  Locally, with respect to complex coordinates
  $z_1, \dots, z_n$
  with $z_j = x_j + i y_j$, a complex
  $k$-form is given by
  \[
    \sum_{\substack{I, J \subseteq \{1, \dots, n\} \\ |I| + |J| = k}} f_{I, J}\,
    dx_{i_1} \wedge \cdots \wedge dx_{i_r}
    \wedge dy_{j_1} \wedge \cdots \wedge dy_{j_s},
  \]
  where the $f_{I, J}$ are smooth
  complex-valued functions. Then we
  have a complex
  \[
    \begin{tikzcd}
      0 \ar[r] & \underline{\C}
      \ar[r, "d"] & \mathcal{A}^0_\C
      \ar[r, "d"] & \mathcal{A}^1_\C
      \ar[r, "d"] & \cdots
    \end{tikzcd}
  \]
  which is a fine resolution
  of $\underline{\C}$. Thus we have
  $H^k(X, \underline{\C}) \cong H^k_{\mathrm{dR}}(X, \C)$.
\end{remark}

\begin{remark}
  Observe that there is a decomposition
  \[
    \mathcal{A}^k_\C
    = \bigoplus_{p + q = k} \mathcal{A}^{p, q},
  \]
  where $\mathcal{A}^{p, q}$ is the sheaf
  of $(p, q)$-forms. Writing
  $z_j = x_j + i y_j$ and $\overline{z}_j = x_j - i y_j$, we can locally write a
  $(p, q)$-form in the form
  \[
    \sum_{|I| = p, |J| = q} f_{I, J}\,
    d z_I \wedge d \overline{z}_J,
  \]
  where the $f_{I, J}$
  are smooth complex-valued functions. Then
    $d : \mathcal{A}^k_\C \to \mathcal{A}^{k + 1}_\C$
  decomposes  as $d = \partial + \overline{\partial}$, with
  \[
    \partial (f_{I, J}\, dz_I \wedge d\overline{z}_J)
    = \sum_{i = 1}^n \frac{\partial}{\partial z_i} f_{I, J} dz_i \wedge dz_I \wedge d\overline{z}_J
  \]
  and $\overline{\partial}$ defined similarly.
  This decomposition satisfies the following properties:
  \begin{enumerate}
    \item $\partial^2 = \overline{\partial}^2 = 0$;
    \item $\partial \overline{\partial} + \overline{\partial} \partial = 0$;
    \item $\partial A^{p, q} \subseteq A^{p + 1, q}$ and $\overline{\partial} A^{p, q} \subseteq A^{p, q + 1}$;
    \item $\OO_X = \ker(\overline{\partial} : \mathcal{A}^{0, 0} \to \mathcal{A}^{0, 1})$
      and $\Omega_X^p = \ker(\overline{\partial} : \mathcal{A}^{p, 0} \to \mathcal{A}^{p, 1})$.\footnote{Here $\OO_X$ denotes the sheaf of holomorphic functions and $\Omega_X^p$ denotes the sheaf of holomorphic $p$-forms. Note that $\Omega_X^p$ coincides with the analytification of the sheaf of algebraic $p$-forms when $X \subseteq \PP^n_\C$ is a projective algebraic variety.}
  \end{enumerate}
  Now the complex
  \[
    \begin{tikzcd}
      0 \ar[r] & \Omega_X^p
      \ar[r] & \mathcal{A}^{p, 0}
      \ar[r, "\overline{\partial}"] & \mathcal{A}^{p, 1}
      \ar[r, "\overline{\partial}"] & \cdots
    \end{tikzcd}
  \]
  satisfies the following:
  \begin{enumerate}
    \item $A^{p, q}$ is fine
      ($\OO_X$ has partitions of unity);
    \item exactness (this follows
      by the antiholomorphic Poincar\'e lemma).
  \end{enumerate}
  Thus we get that
  \[
    H^q(X, \Omega_X^p)
    = H^q(\Gamma(X, \mathcal{A}^{p, \bullet}))
    = \frac{\ker(\overline{\partial} : \mathcal{A}^{p, q}(X) \to \mathcal{A}^{p, q + 1}(X))}{\im(\overline{\partial} : \mathcal{A}^{p, q - 1}(X) \to \mathcal{A}^{p, q}(X))}
    =: H^{p, q}.
  \]
  We call the $H^{p, q}$ the
  \emph{Dolbeault cohomology} of $X$.
  Then the Hodge decomposition theorem
  says that
  \[
    H^k_{\mathrm{sing}}(X, \C)
    = \bigoplus_{p + q = k} H^{p, q}.
  \]
\end{remark}
